\section{Hardware-enforced security barriers}

\subsection{Privilege levels}

\begin{frame}{Kernel/Userland isolation}
  \begin{itemize}
    \item We expect some security guarantees from modern systems:
    \begin{itemize}
      %% This is of course enforced by the kernel, and will be defeated if
      %% a process can just access the entire RAM, or even disk when swapping.
      \item Process isolation (resources, address space)
      \item User isolation
      %% This will not work if processes have access to the disk's HW.
      \item Access control on files
    \end{itemize}
    %% This entity will need 
    \item We need an entity in charge of enforcing those guarantees
    \item That entity is the operating system's kernel, this is one of its main roles
  \end{itemize}
\end{frame}


\begin{frame}{Kernel/Userland isolation}
  \begin{itemize}
    \item The kernel has \textbf{privileged} access to the hardware
    \item On older CPUs, or some architectures, it even has \textbf{unfiltered} access
    \item Needs to be isolated from useland even though both run on the same hardware
    \item The key is the \textbf{privilege level}
    \begin{itemize}
      \item x86 defines 4 "protection rings", tracked in 2 bits of the Code Segment Selector register
            Only 2 are actually used: ring 0 for the kernel, ring 3 for userland
      \item ARM has 4 "Exception Levels" (EL), with the kernel in EL1 and userland in EL0.
            The current EL is tracked in the \href{https://developer.arm.com/documentation/ddi0595/2021-06/AArch64-Registers/CurrentEL--Current-Exception-Level}{CurrentEL}
 register
    \end{itemize}
  \item The privilege level is checked e.g. when accessing memory
  \item Some instructions can only be executed from eleveated privilege levels
  \end{itemize}
\end{frame}

\begin{frame}{Protection rings}
\includegraphics[width=0.5\textwidth]{slides/security-hw-sec-barriers/rings.pdf}  
\end{frame}

\begin{frame}[fragile] 
  \frametitle{System Calls (1/2)}
  \begin{itemize}
    \item A system call allows the user space to request services from the
          kernel by executing a special instruction that will switch to the
          kernel mode (\manpage{syscall}{2})
    \begin{itemize}
      \item When executing functions provided by the libc (\code{read()},
            \code{write()}, etc), they often end up executing a system call.
    \end{itemize}
    \item System calls are identified by a numeric identifier that is passed
          via the registers.
    \begin{itemize}
      \item The kernel exports some defines (in \code{unistd.h}) that are named
            \code{__NR_<sycall>} and defines the syscall identifiers.
    \end{itemize}
  \end{itemize}
  \begin{block}{}
    \begin{minted}[fontsize=\small]{C}
#define __NR_read 63
#define __NR_write 64 
    \end{minted}
  \end{block}
\end{frame}
    
\begin{frame}[fragile]
  \frametitle{System Calls (2/2)}
  \begin{itemize}
    \item The kernel holds a table of function pointers which matches these
          identifiers and will invoke the correct handler after checking the
          validity of the syscall.
    \item System call parameters are passed via registers (up to 6).
    \item When executing this instruction the CPU will change its execution
    state and switch to the kernel mode.
    \item Each architecture uses a specific hardware mechanism
    (\manpage{syscall}{2}) 
  \end{itemize}
  \begin{block}{}
    \begin{minted}[fontsize=\small]{asm}
    mov w8, #__NR_getpid
    svc #0
    tstne x0, x1
    \end{minted}
  \end{block}
\end{frame}


\begin{frame}{Virtual address translation}
  \begin{itemize}
    \item When a core is running a user program it is not running the kernel
    %% Running memory accesses through the kernel would be so painfully slow, the performance
    %% penanlty is not acceptable to even hypervisors
    \item The hardware itself must be able to enforce memory protection
    \begin{itemize}
      \item this is where the MMU comes in
    \end{itemize}
    %% ARMv8 actually calls pages "blocks", and refers to the different page sizes as "granules"
    \item The CPU accesses \textbf{virtual}, not \textbf{physical} addresses
    %% Note: would be cool to have the logic checking the physical address mapping
    \item Each process has a different virtual address space
    \item Virtual memory is organized into a hierarchy of tables, the lowest level being pages.
          see \kdochtml{mm/page_tables}
    \item Table elements are descriptors, and contain several attributes
    \item Documentation:
    \begin{itemize}
      \item ARMv8-A: \href{https://developer.arm.com/documentation/100940/0101/}{Armv8-A Address Translation}
      \item x86: \href{https://www.intel.com/content/www/us/en/developer/articles/technical/intel-sdm.html}{Intel Software Developer's Manual, Volume 3a, Chapter 5}
    \end{itemize}
  \end{itemize}
\end{frame}

\begin{frame}{Memory attributes}
  \begin{itemize}
    \item Page descriptors use a hardware-specific format:

    \item \kfunc{PTE_USER} (ARM64) / \kfunc{_PAGE_BIT_USER} (x86) indicates userland
    \item An entry can be marked non-executable:
    \begin{itemize}
      \item on x86, using the NX bit (\kfunc{_PAGE_BIT_NX})
      \item on 64-bit ARM, split between EL0 (\kfunc{PTE_UXN}) and EL1+ (\kfunc{PTE_PXN})
    \end{itemize}
    %% Userland (ring 3/EL0) should not be able to execute kernel (ring 0/EL1) code
    \item Userland (ring 3/EL0) execution from kernel (ring 0 /EL1) is conditional:
    \begin{itemize}
      \item x86: depends on SMEP, see \href{https://www.intel.com/content/www/us/en/developer/articles/technical/intel-sdm.html}{Intel SDM, Vol 3a, 5.6.1}
      \item 64-bit ARM: this is why \textbf{XN} is split into \kfunc{PTE_UXN} and \kfunc{PTE_PXN}
    \end{itemize}
    \item On x86, userland (ring 3) access from kernel (ring 0) can be disabled using SMAP
  \end{itemize}
\end{frame}

\begin{frame}{Memory attributes (ARMv8)}
\includegraphics[width=0.6\textwidth]{slides/security-hw-sec-barriers/armv8_descriptor.pdf}  
  \begin{itemize}
    \item NS set to 1 means we are in normal world
    \item bit 7 (AP.1) encodes read-only
    \item bit 6 (AP.0) encodes unprivileged (EL0) access
    \newline
    \begin{tabular}{| l || c | c | c | c | c |}
      \hline
      AP & bit 7 & bit 6 & EL0 & EL1+ \\
      \hline
      00 & 0 & 0 & No access & RW \\
      \hline
      01 & 0 & 1 & RW & RW \\
      \hline
      10 & 1 & 0 & No access & RO \\
      \hline
      11 & 1 & 1 & RO & RO \\
      \hline
    \end{tabular}
  \end{itemize}
\end{frame}

\begin{frame}[fragile]
  \frametitle{Experimenting with memory attributes}
  %% ptr will be 0x7XXXXXXXX000 because with NULL, mmap() will return
  %% a page-aligned address (we are supposing a 4K page, so PAGE_SHIFT=12)
  %% We know that it will be 12 char (6 bytes) because 64-bit architectures
  %% mostly use 48-bit physical addresses, and the MMU enforces it on
  %% virtual addresses as well.
  %% AArch64 is more complicated, because it supports physical addresses
  %% from 40 bytes up to 56 bytes (in ARMv9.3).
  %% However, in all cases, the MSB is used as a differentiator between
  %% userland (0) and kernel (1), and the adress is sign-extended.
  %% An x86-64 address with bit 47 at 0 and any of bits 63-48 at 1 is
  %% non-canonical and will be rejected by the MMU.
  %% Acces to 4097 will succeed, because it is located past the first
  %% 4096-byte page that had its permissions changed, whereas access to
  %% 1025 will SEGFAULT, because although it is past the 1024 mmap range,
  %% it is still located on the same page, and therefore had its protection
  %% modified. Of course, we know that ptr is the start of a page because
  %% that is what mmap does when passed NULL as 1st argument.
  \begin{minted}[fontsize=\scriptsize]{c}
#include <stdio.h>
#include <stdlib.h>
#include <sys/mman.h>

int main(int argc, char *argv[])
{
        unsigned char *ptr = mmap(
            NULL, 4096 * 2, PROT_READ | PROT_WRITE, MAP_PRIVATE | MAP_ANONYMOUS, -1, 0
        );

        printf("ptr: %p\n", ptr);                               // 0x7XXXXXXXX000
        ptr[0] = 0xFF;
        printf("ptr[1025]: 0x%02x\n", ptr[1025]);               // 0x00
        printf("ptr[4097]: 0x%02x\n", ptr[4097]);               // 0x00

        unsigned int ret = mprotect(ptr, 1024,  PROT_NONE);

        printf("ptr[4097]: 0x%02x\n", ptr[4097]);               // 0x00
        printf("ptr[1025]: 0x%02x\n", ptr[1025]);               // SIGSEGV
}
  \end{minted}
\end{frame}

\begin{frame}{Limits to this model}
  \begin{itemize}
    \item Although this model is widespread, it has limitations
    \item Sometimes, the kernel is untrusted:
    \begin{itemize}
      \item Host virtual machine kernel
      \item "Custom ROMs"
    \end{itemize}
    \item Or it is not trusted \textbf{enough}:
    \begin{itemize}
      \item Large attack surface
      \item Not Invented Here
      \item Money (not just the user's) is at stake
    \end{itemize}
    \item This gave rise to multiple hardware isolation technologies
  \end{itemize}
\end{frame}

\subsection{ARM TrustZone}

\begin{frame}{TrustZone / Secure World}
  \begin{itemize}
    \item \href{https://developer.arm.com/documentation/100935/0100/The-TrustZone-hardware-architecture-?lang=en}{TrustZone} is an ARM-only hardware isolation mechanism  
    \item It is implemented on the \href{https://developer.arm.com/Architectures/A-Profile\%20Architecture}{Application(A)-profiles} of ARMv7-9, this is what we discuss
    \item Also implemented on the ARMv7-M (Microcontroller) profile, but differently
    \item Splits the system into 2 \textbf{worlds}: Secure and non-Secure
    \item They are also called \textbf{Trusted} / \textbf{Rich} Execution Environment (TEE/REE)
    \item This is \textbf{orthogonal} to Exception Levels:
    \begin{itemize}
      %% Technically, Chapter D1.1 the Arm Architecture Reference Manual for A profile Architecture
      %% ()
      %% lists both EL2 and EL3 as optional unless on ARMv9 with the Realms Management Extenstion.
      %% In practice, buying an A profile core without EL3 is probably not worth it.
      \item EL3 (Secure Monitor) only exists in \textbf{Secure} world
      \item EL2, if it exists, can exist in \textbf{Secure}, \textbf{Non-Secure}, or both
      %% This is not accurate for ARMv9 with RME (Realms Mangement Extensions)
      \item EL1 and EL0 must exist in \textbf{both} security states
    \end{itemize}
  \end{itemize}
\end{frame}

\begin{frame}{TrustZone / Secure World}
  \center\includegraphics[width=0.7\textwidth]{slides/security-hw-sec-barriers/Exception_Levels.pdf}
\end{frame}

\begin{frame}{World transitions: Secure to Non-Secure}
  \begin{itemize}
    \item The CPU starts executing at the highest EL: EL3, so in Secure World
    \item It then sets up the rest of Secure World, then Normal World, in stages
    \item Similar to the EL, the CPU keeps track of the security state in a register
    \begin{itemize}
      \item SCR is the \textit{Secure Configuration Register (SCR)}, bit 0 is the \textit{NS} bit
      \item SCR can only be accessed in Secure EL1-3
      \item Going from Secure World back to Non-Secure is simply setting \textit{NS} to 1
      \item The Security model expects only the Secure Monitor to do this transition
    \end{itemize}
    \item The ARM spec leaves non-EL3 access to SCR as \textit{undefined}
  \end{itemize}
\end{frame}

\begin{frame}{World transitions: Back to Secure}
  \begin{itemize}
    \item The systems transitions from Non-Secure to Secure via exceptions, either:
    \begin{itemize}
      \item asynchronous, via some \textit{FIQ} interrupts, depending on \textit{SCR.FIQ}
      \item synchronous, generated by an \textit{smc} (\textit{secure monitor call}) instruction
    \end{itemize}
    \item \textit{smc} cannot be called from EL0, so userland would first need to issue a \textit{svc}
    \item The authority is the \textit{Secure Configuration Register (SCR)}'s NS bit
    \item Accessing SCR from non-secure or from S-EL0 will fault.
    \item Only the Secure Monitor (running at EL3) \textbf{should} change it.
  \end{itemize}
\end{frame}

\begin{frame}{\href{https://www.trustedfirmware.org/}{Arm Trusted Firmware}}
  \begin{itemize}
    \item Trusted Firmware is an open-source reference Secure World implementation
    \item It is the default Secure World for most ARM platforms
    \item Comes in two flavours:
    \begin{itemize}
      \item Trusted Firmware A (for the Application profile), or
            \href{https://www.trustedfirmware.org/projects/tf-a}{TF-A}
      \item Trusted Firmware M (for the Microcontroler profile)
            \href{https://www.trustedfirmware.org/projects/tf-m}{TF-M}
    \end{itemize}
    \item There is now a Rust implementation of TF-A called Rusted Firmware A
          (\href{https://www.trustedfirmware.org/projects/rf-a}{RF-A})
    \item Most importantly, TF-A implements the Secure Monitor
    \item Uses a specific nomenclature to designate boot components
  \end{itemize}
\end{frame}

\begin{frame}{Trusted Firmware A (TF-A)}
  \begin{itemize}
    \item TF-A is split into 5 steps:
    \begin{itemize}
      \item BootLoader 1, or BL1: the SoC vendor's ROM code
      \item BootLoader 2, or BL2: "Trusted Boot firmware"
      \item BootLoader 3-1, or \href{https://elixir.bootlin.com/arm-trusted-firmware/v2.14.0/source/bl31}{BL31}:
            EL3 Runtime
      \item BootLoader 3-2, or \href{https://elixir.bootlin.com/arm-trusted-firmware/v2.14.0/source/bl32}{BL32} (optional):
            Secure-EL1 payload (Trusted OS, Trusted kernel, TEE)
      \item BootLoader 3-3, or BL33: Non-trusted firmware, the "normal" bootloader
    \end{itemize}
    \item Only BL31 and BL32 are resident after boot is complete:
    \begin{itemize}
      \item BL31 is responsible for routing Secure exceptions (and smcs) to the Secure World
      \item BL32 handles some of these exceptions and smcs, potentially on behalf of TAs
    \end{itemize}
  \end{itemize}
\end{frame}

\begin{frame}{ARMv8 boot sequence: TF-A names}
\includegraphics[width=0.7\textwidth]{common/armv8-boot-sequence-generic.pdf}
\end{frame}

\begin{frame}{Trusted Firmware A (TF-A)}
  \begin{itemize}
    \item In a "standard" boot, where U-Boot is used as boot loader:
    \begin{itemize}
      \item BL1: the vendor bootROM, non-resident
      \item BL2: is implemented by the \href{https://docs.u-boot.org/en/latest/usage/spl_boot.html}{Second Program Loader} (SPL)
      \item BL31: the runtime component of TF-A resident in EL3
      \item BL32: OP-TEE
      \item BL33: U-Boot proper
    \end{itemize}
  \end{itemize}
\end{frame}

\begin{frame}{ARMv8 boot sequence}
\includegraphics[width=0.7\textwidth]{common/armv8-boot-sequence-common-names.pdf}
\end{frame}

\begin{frame}{Rationale}
  \begin{itemize}
    \item The aim of TrustZone is to create a Trusted Execution Environment (TEE)
    \item The isolation is backed by the hardware
    \item Some hardware resources are only accessible to the TEE
    \item Sounds familiar, why is it not simply the kernel?
    \begin{itemize}
      \item attack surface: the TEE should be as small as possible
      \item usage: can you ensure attackers do not get \textit{root}
      \item virtualization: offer services securely to kernels you don't trust
    \end{itemize}
    \item The key concept is defense in depth
  \end{itemize}
\end{frame}

\begin{frame}{Trusted Execution Environments}
  \begin{itemize}
    \item The TEE is the part of Secure World that does things for the user
    \item This is an A profile concept, the M profile equivalent is Secure Processing Environment (SPE)
    \item There are different implementations of TEEs:
    \begin{itemize}
      \item Trustsonic's \href{https://www.trustonic.com/secure-os/technology/}{Kinibi}
      \item Samsung's \href{https://www.samsungknox.com/en}{Knox}
      \item Qualcomm's \href{https://docs.qualcomm.com/doc/80-88500-4/topic/78_Qualcomm_TEE_5_0.html}{QTEE}
      \item Huawei's iTrustee
    \end{itemize}
    \item Some are even Open Source:
    \begin{itemize}
      \item Android's \href{https://source.android.com/docs/security/features/trusty}{Trusty},
            based on Little Kernel (LK)
      \item Nvidia's Trusted Little Kernel (TLK), also based on LK: \code{git://nv-tegra.nvidia.com/3rdparty/ote_partner/tlk.git}
      \item OP-TEE, the reference Open Source TEE
    \end{itemize}
  \end{itemize}
\end{frame}

\begin{frame}{Trusted Execution Environments}
  \begin{itemize}
    \item The TEE is an optional component
    \item It does not normally offer service to Normal World directly
    \item It does provide HW access to Trusted Applications
    \item The TEE is the Secure World equivalent of the kernel, running in S-EL1
    \item For further isolation, ARMv8 introduces Secure Partitions
    \begin{itemize}
      \item Essentially run different instances of TEEs in isolated HW
      \item A Secure Partition Manager runs at S-EL2 (SPMC) and (S-)EL3 (SPMD)
    \end{itemize}
    \item ARMv9 pushed this even further with the Realms extension
  \end{itemize}
\end{frame}

\begin{frame}{Trusted Applications}
  \begin{itemize}
    \item Trusted Applications are the secure world "userland"
    \item OP-TEE has 2 main types of Trusted Apps:
    \begin{itemize}
      \item Pseudo-TAs, statically compiled into the OP-TEE core binary
      \item Usermode TAs are dynamically loaded by OP-TEE and run at S-EL0.\\
    \end{itemize}
    \item Usermode TAs can be further subdivided into:
    \begin{itemize}
      \item Early TAs: still stored in OP-TEE core, in a \code{data} section
      \item REE FS TAs: loaded by the REE via \code{tee-supplicant}.
    \end{itemize}
    \item Because REE FS TAs are loaded from the REE, the TEE needs to authenticate them
    \begin{itemize}
      \item OP-TEE expects REE FS TAs to be \href{}{signed}
      \item the OP-TEE core embarks the public part of
            \href{https://elixir.bootlin.com/op-tee/4.9.0/source/keys/default_ta.pem}{\code{default_ta.pem}}
      \item do \textbf{NOT} ship OP-TEE with that key
    \end{itemize}
  \end{itemize}
\end{frame}

\begin{frame}{Usecases}
  \begin{itemize}
    \item "Software HSM": this is why OP-TEE has a \href{https://github.com/OP-TEE/optee_os/tree/master/ta/pkcs11}{PKCS\#11 TA}
    \item \href{https://github.com/OP-TEE/optee_ftpm}{fTPM}: OP-TEE 
    \item Payment processing
    \item DRM
    \item Biometrics (fingerprint unlock)
    \item Hardware-isolated detection of non-secure malware
  \end{itemize}
\end{frame}

\begin{frame}{Issues}
  \begin{itemize}
    \item TrustZone is, ultimately, another security layer
    \item The security guarantees it offers are only as strong as the implementation
    \item Although it is hardware-backed, it is only useful if software is involved
    \item Like all software, Secure World software can have bugs
    \begin{itemize}
      \item \href{https://blog.thalium.re/posts/pivoting_to_the_secure_world/}{Arbitrary code execution in a Samsung TEE}
      \item \href{https://www.keysight.com/blogs/en/authors/federico-menarini}{Other TEEGris vulnerabilities}
      \item \href{https://blog.quarkslab.com/attacking-the-arms-trustzone.html}{S-EL0 stack buffer overflow in Qualcomm's TEE}
      \item \href{https://blog.quarkslab.com/a-deep-dive-into-samsungs-trustzone-part-3.html}{EL3 execution via Samsung Kinibi}
    \end{itemize}
    \item Secure world configuration can be very HW-dependent
  \end{itemize}
\end{frame}


